\section{The Conversion of Valentin Tomberg}

Much ink (virtual and real) has been spilled regarding the conversion of \textbf{Valentin Tomberg} from Anthroposophy to
Roman Catholicism. Yet, as we saw in the conversion of \textbf{Rene Guenon}, such a move cannot be understood in the
conventional sense as the rejection of one thing and the adoption of another.

Nevertheless, there are those who are convinced that Tomberg rejected so-called New Age teachings to become a Catholic,
and therefore attempt to follow him in that path. As a matter of fact, it seems that the Internet is replete with
converts (or reverts) who are quite enthusiastic in promoting their new-found faith, usually to excess. While we think
such conversions are a good thing in general, that is not at all Tomberg's message.

In Letter XI, Tomberg explain his reasons for entering the Church.

\begin{quotationx}
The way of Hermeticism, solitary and intimate as it is, comprises authentic experiences from which it follows that the
Roman Catholic Church is, in fact, a depository of Christian spiritual truth, and the more one advances on the way of
free research for this truth, the more one approaches the Church. Sooner or later one inevitably experiences that
spiritual reality corresponds — with an astonishing exactitude — to what the Church teaches. 

\end{quotationx}
He then lists several specific teachings:

\begin{itemize}
\item There are guardian angels 
\item There are saints who participate actively in our lives 
\item The Blessed Virgin is real, as she is understood, worshiped, and portrayed 
\item The sacraments are effective and there are seven of them 
\item The three sacred vows of obedience, chastity, and poverty constitute the very essence of all authentic
spirituality 
\item Prayer is a powerful means of charity 
\item The ecclesiastical hierarchy reflects the celestial hierarchical order 
\item The Holy See and the papacy represent a mystery of divine magic 
\item Hell, purgatory, and heaven are realities 
\item The Master himself abides with his Church 
\item The Master is always findable and meetable there 
\end{itemize}
\paragraph{The Common Believer}
Any cursory reading of the Meditations shows that it is replete with references to ideas, systems, books, and people
that are certainly precursors to the New Age teachings of today. Yet Tomberg did not reject them in toto, not at all.
To the contrary, he makes the remarkable claim that following Hermetic teaching to its depths led to his “conversion”.
The reasons for this must be explored in what follows. But first of all, note that Tomberg does not claim any sort of
superiority; rather, he acknowledges a solidarity with common believers, as expressed in Letter IV.

\begin{quotationx}
For the Hermetic-philosophical sense has more in common with the plain and sincere faith of simple people than abstract
metaphysics has.

\begin{itemize}
\item
For the common believer, God lives; likewise for the Hermeticist. 
\item 
The believer addresses himself to saints and Angels; for the Hermeticist they are real. 
\item
The believer believes in miracles; the Hermeticist lives in the presence of miracles. 
\item
The believer prays for the living and the dead; the Hermeticist dedicates all his efforts in the domain of sacred magic
to the good of the living and the dead. 
\item
The believer esteems all that which is traditional; the Hermeticist does likewise. 
\end{itemize}
\end{quotationx}

\paragraph{Rejection of Alternatives}
In the Introduction to \textit{Inner Development}, we are informed that Tomberg initially tried to align himself with the
Christian Community and then with Russian Orthodoxy.

The Christian Community was formed by some of Rudolf Steiner's followers; they have no dogmas, although they
have priests and seven sacraments. However, dogmas are not an affront to free will as the Community claims, but rather
they are living symbols of a higher spiritual reality as Tomberg came to realize. Typically, organisations that reject
dogmas tend to converge to liberalism.

Tomberg attempted to work with the Christian Community by introducing a cult of Mary-Sophia. Emil Bock reportedly said
to him: “We have Michael, that's enough! We don't need Mary-Sophia.”

As we will see in the next section, that is absolutely contrary to Tomberg's purpose and mission.

Orthodoxy lacks a complete hierarchy, in particular, the papacy which Tomberg regards as a mystery of divine magic.
Russian Orthodoxy retains a notion of being the Third Rome, with an Emperor and the Patriarch as Pope. From the
esoteric perspective, this is a caricature of the true teaching, whose real source is suspect. Most of the early popes
came from the East; they were more aware of the celestial hierarchy and divine magic. Someday, there will be a Russian
pope ruling in the first Rome.

\paragraph{The Blessed Virgin}
Tomberg insists that the Blessed Virgin is real, as she is understood, worshipped, and portrayed in the Church. This
means that he accepts the four Marian dogmas, not as beliefs but as a personal experience. These are:

\begin{itemize}
\item Mary as the Mother of God 
\item Perpetual Virginity 
\item The Immaculate Conception 
\item The Assumption 
\end{itemize}
Although it is not so common today, when I was a schoolboy, we would attend yearly novena to the Blessed Virgin. We
would be given scapulars or medals. Although I might not have fully understood their significance at the time, I have
always lived in the security of Mary's promises of protection. Tomberg describes the esoteric meaning of
these promises:

\begin{quotationx}
Every Hermeticist who truly seeks authentic spiritual reality will sooner or later meet the Blessed Virgin. This meeting
signifies, apart from the illumination and consolation that it comprises, protection against a very serious spiritual
danger. For he who advances in the sense of depth and height in the “domain of the invisible” one day arrives at the
sphere known by esotericists as the “sphere of mirages” or the “zone of illusion”. This zone surrounds the earth as a
belt of illusory mirages. It is this zone which the prophets and the Apocalypse designate “Babylon”. The soul and the
queen of this zone is in fact Babylon, the great prostitute, who is the adversary of the Virgin. … One cannot traverse
it without the protection of the “mantle of the Blessed Virgin”. 

\end{quotationx}
Mary is celebrated as the Queen of Heaven. Tomberg has a deeper insight that takes that teaching much further. In Letter
XI, he asserts:

\begin{quotationx}
The day when it is achieved will be the day of a new festival — the festival of the coronation of the Virgin
on earth. For then the principle of opposition will be replaced on earth by that of collaboration. This will be the
triumph of life over electricity. And cerebral intellectuality will then bow before Wisdom (SOPHIA) and will unite with
her. 

\end{quotationx}
The Virgin will be not only Queen of Heaven, but also Queen of the Earth. Of course, this is confirmed in the
\emph{Dogmatic Constitution}, which calls Mary the Queen of the Universe.

The triumph of life over electricity refers to the Hermetic teaching of electricity, or electro-magnetism, arising from
lower forces. Just as a reminder, Tomberg explains:

\begin{quotationx}
The fruit of the Tree of Knowledge of Good and Evil — the fruit of the polarity of opposites — is therefore electricity; and electricity entails fatigue, exhaustion. death. Death is the price that is paid for the
knowledge of good and evil, i.e. the price of life amidst opposites. For it is electricity — physical,
psychic and mental — which was introduced into the being of Adam-Eve. and thereby into the whole of
life-endowed Nature, from the moment that Adam-Eve entered into communion with the tree of opposites, that is to say
with the principle of electricity. And it is thus that death entered into the domain of life-endowed Nature. 

\end{quotationx}
Therefore, the next dogma will be Mary as Co-Redemptrix.

\paragraph{Way, Truth, and Life}
Tomberg rejected Anthroposophy (although not Steiner himself), describing it as

\begin{quotationx}
a movement for cultural reform (art, education, medicine, agriculture) deprived of living esotericism, i.e. without
mysticism, without gnosis and without magic, which have been replaced by lectures, study and intellectual work aiming
at establishing a concordance between the writings and stenographed lectures of the master. 

\end{quotationx}
In other words, it has Truth but not the Life. It is locked in concepts, meaning that one learns the concepts first,
then tries to have the experience, whereas it should be vice versa. Tomberg explained this in a letter to Bernhard
Martin:

\begin{quotationx}
First they [i.e., Anthroposophists] have a world of formulated concepts and then try to arrive at experience. But the
concepts hold them shut within their world: the spiritual world remains silent, because they are the ones talking about
the spiritual world; they don't let it speak. It's otherwise with people [like Jung]; in
silence they let the spiritual world speak. And the spiritual world speaks in symbols — i.e. in mystery
speech — today just like before. 

\end{quotationx}
In other words, it is necessary to treat the concepts as symbols, as the symbol is understood in the Meditations. It is
an invitation to a personal meditation, not a univocal concept to be learned. In \emph{Covenant of the Heart}, Tomberg
is more explicit:

\begin{quotex}
Alas it happened, however, for reasons which we need not go into here, that Rudolf Steiner gave his work the form of a
science, so-called “spiritual science”. Thereby the third aspect of the indivisible threefoldness of the Way, the
Truth, and the Life was not given enough attention. For the scientific form into which the logic of the Logos had to be
cast, and by which it was limited, left little room for pure mysticism and spiritual magic, that is, for Life. So there
is in Anthroposophy a magnificent achievement of thought and will — which is, however, unmystical and
unmagical, i.e. in want of Life. Rudolf Steiner himself was conscious of this essential lack. Therefore, it was with a
certain amount of hope that he indicated the necessary appearance of a successor (the Bodhisattva), who would remedy
this lack and would bring the trinity of the Way, the Truth, and the Life to full fruition. 

\end{quotex}
\paragraph{Knowledge as Intuition}
Tomberg is trying to get us out of our heads in order to experience a higher type of knowledge. In \emph{Covenant of the
Heart}, he explains:

\begin{quotationx}
truth is based on “intuition [which] is not attained through practical knowledge or intellectual consideration
(reflection), but through direct experience of reality … `an evolving revelation from the inner being of
man’ ... and `a direct grasping of the being of things'.

\end{quotationx}
He then goes on:

\begin{quotationx}
For those who experience it, this form of knowledge counts as the highest because it is experienced … as the result of
the most profound contemplation and the greatest concentration, in comparison with which that of intellectual
consideration and the practical knowledge gained by way of observation appears superficial. However, it does not count
in the slightest way as knowledge (let alone as the highest form of knowledge) for the scientific disciplines which, as
such, lay claim to being of general validity. For the scientific approach is not to strive simply for the truth, but
rather to strive for that brand of truth which is of general validity, i.e. that which can be comprehended
fundamentally by everyone bestowed with healthy understanding and faculties of perception, and which should thus be
concurred with. A scientific discipline, whether a spiritual-scientific or a natural-scientific discipline, does not
want to, and is not able to, address itself only to those people who are capable of the concentration and inner
deepening necessary for intuition. Were it to do so, it would then not be scientific, i.e. generally comprehensible and
provable. Rather, it would be “esoteric”, i.e. a matter for an elite group of special people. In this sense theology is
also “science” since, assuming the authority of Scripture and the Church are acknowledged, it can be comprehended and
tested by all believers. 

\end{quotationx}
Direct spiritual knowledge achieved through intuition is \emph{personal}—never general or universal, i.e.
scientific in the conventional sense. That is why Hermetism is not one philosophical system among many systems, nor one
scientific theory among other competing views, nor the foundation of a new religion. In other words, it is not
expansive in the horizontal sense, but rather a matter of depth, i.e., a deepening of understanding.

\paragraph{Postscript}
So, to simplify, we can define:

\begin{itemize}
\item The Truth as the understanding of concepts 
\item The Way as the deepening of that understanding 
\item The Life as direct intuition 
\end{itemize}
An example might be this:

\begin{itemize}
\item Hegel's system of Absolute Idealism, or perhaps other similar systems 
\item Steiner created not just a thought system, but also proposed a path of spirit and soul development 
\item Tomberg opens up the meaning of symbols and intuition 
\end{itemize}

\flright{\itshape Posted on 2018-08-17 by Cologero}
